\documentclass[12pt]{article}
\usepackage[margin=1in]{geometry}
\usepackage[T1]{fontenc}
\usepackage[utf8]{inputenc}
\usepackage{setspace}
\usepackage{parskip}

\title{The Last Remaining Asymmetry:\\Fairness, Competition, and the Future of Fertility}
\author{Kangbin Yim\thanks{Email: yim@sch.ac.kr (Prof., SCH University, Korea)}}
\date{May 2026}

\begin{document}
\maketitle

For decades, the decline in fertility has been explained almost entirely in economic terms. Housing is too expensive, education costs too much, employment is unstable, and the opportunity cost of raising children keeps rising. Governments have responded accordingly, with subsidies, tax incentives, childcare programs, and an expanding architecture of family support. And yet, across many of the world's most developed societies, birth rates continue to fall. The persistence of the trend, in the face of so much spending designed to reverse it, invites a difficult possibility: that the problem may not be purely economic, and that something deeper is at work beneath the familiar surface of cost.

Most discussions of fertility concentrate on cost. Far fewer consider the structure of civilization itself. Yet civilization has moved, with remarkable consistency, in a single direction—it has steadily reduced asymmetry. The distance between rulers and the ruled has narrowed. The line dividing aristocrat from commoner has faded. Education has become more accessible, political participation more inclusive, and economic life less dependent on inherited privilege. Seen this way, much of human history can be read as one long expansion of fairness.

But fairness carries a consequence that is easy to overlook. It does not abolish competition; it universalizes it. In a rigidly stratified society, most people are excluded from the contest before it begins. In a more symmetrical one, far more people step into the same arena. The result is not less competition but more of it—broader, denser, and shared. The modern world, then, is not simply fairer than what came before. It is also more competitive, precisely because it is fairer.

Nowhere is this clearer than in the changing relationship between men and women. For most of human history, the sexes occupied substantially different social roles, and the reasons were not only cultural but functional. Survival leaned heavily on physical strength, on exposure to environmental risk, on warfare, hunting, and labor-intensive production. Men and women entered society through different pathways because society itself was organized around different demands. Whether those arrangements were just is a separate question; what matters here is that the competition was largely separated. Men and women were often contending within different domains rather than against one another.

Modern civilization has dismantled those divisions piece by piece. Industrialization diminished the premium on physical strength. Mass education widened access to knowledge. Digital technology lowered the barriers to information. Professional opportunity, economic independence, and political rights all expanded. As they did, men and women increasingly came to participate within the same systems of competition. They attend the same schools and compete for the same university places, the same jobs, the same promotions, the same positions of leadership, and the same definitions of success. The direction is unmistakable: modern society is not only producing greater equality but assembling a single, shared competitive arena.

Artificial intelligence and robotics are now accelerating this convergence. Machines have already absorbed much of the physical labor that once favored male strength, and algorithms are beginning to automate the cognitive labor that once created other forms of differentiation. Robots take on the physically demanding tasks; software takes on the analytical ones. As these technologies advance, biological sex becomes progressively less relevant to economic productivity and professional achievement. The long-term implication is striking. For the first time in human history, men and women may come to compete within almost entirely symmetrical social structures.

And it is exactly here, at the threshold of near-complete symmetry, that a paradox surfaces. As asymmetries vanish one after another, a single asymmetry grows more conspicuous: childbirth. Pregnancy occurs within the female body; so does birth itself. No political reform can erase this, no economic policy can redistribute it, no legal framework can fully equalize it. It remains a biological fact. And paradoxically, the more symmetrical a society becomes, the more starkly that fact stands out.

In earlier societies, asymmetry was everywhere. Men frequently bore greater physical risk through war, hunting, and dangerous labor; women bore the burdens of pregnancy, childbirth, and child-rearing. Because society rested on many such asymmetries at once, no single one dominated the collective imagination. Today most of them have weakened—educational, political, professional, economic. As they recede, childbirth is increasingly left standing alone, one of the few asymmetries that cannot be removed from the system.

The significance of this reaches beyond biology. In a society organized around fair competition, childbirth may come to be experienced as an asymmetric burden within an otherwise symmetrical arena. This is not to say that women consciously refuse childbirth because they judge it unfair; human decisions are far more tangled than that. Economic conditions matter, as do cultural values, personal aspirations, and public policy. The claim is subtler. It is that the meaning of childbirth may shift when set against a background of expanding fairness and intensifying competition. An asymmetry that once stood among many now stands among few; an asymmetry that once lay within separate social pathways now lies within shared ones. For both reasons it becomes more visible—and what grows more visible may also grow more consequential.

This may help explain why fertility decline tends to be most severe in the most developed societies. Such societies are not merely wealthier; they are more symmetrical and more competitive. They offer greater educational opportunity, greater professional mobility, greater individual autonomy—and, increasingly, they place men and women within the same systems of evaluation. The very forces that enlarge freedom and fairness may, at the same time, sharpen awareness of biological asymmetry.

From this vantage point, falling fertility looks less like a simple demographic event and more like a tension between two historical trajectories: civilization's long movement toward fairness and broader competition, and the biological structure of human reproduction. For most of history the two coexisted quietly, cushioned by the many asymmetries society contained. As civilization removes those asymmetries, reproduction itself becomes increasingly exceptional.

The question facing advanced societies may therefore differ from the one policymakers usually pose. It may not be enough to ask how to lower the economic cost of having children. The deeper question may be this: what happens when a society built on fairness and universal competition encounters a biological asymmetry that cannot be fully removed? Artificial intelligence and robotics only sharpen the issue. As technology continues to dissolve distinctions in labor, capability, and social participation, the visibility of biological asymmetry may keep rising. The future of fertility, in the end, may turn less on economics alone than on how civilization chooses to answer its last remaining asymmetry.

\end{document}
